\documentclass{beamer}

\usepackage{graphicx}
\usepackage{color}
\usepackage{listings}
\usepackage{parskip}
\usepackage{hyperref}
\usepackage{pgfplots}
\usepackage{pdfpages}

\frenchspacing

\definecolor{commentcolor}{rgb}{0, 0.7, 0.3}
\definecolor{keywordcolor}{rgb}{0.6, 0, 0.5}
\definecolor{stringcolor}{rgb}{0, 0, 0.9}
\definecolor{logocolor}{rgb}{0.3, 0.3, 0.7}

\lstset {%
	language=[LaTeX]TeX,
	basicstyle=\footnotesize\ttfamily,
	breaklines=true,
	numbers=left,
	tabsize=4,
	commentstyle=\color{commentcolor},
	keywordstyle=\color{keywordcolor},
	stringstyle=\color{stringcolor},
	columns=flexible,
	morekeywords={\therefore, \because, \pgfplotsset, \includepdf, \includegraphics, \maketitle, \tableofcontents},
}

\pgfplotsset{every axis/.append style={
	axis x line=middle, % Make the x-axis position be the middle of the graph
	axis y line=middle, % Make the y-axis position be the middle of the graph
	axis line style={->}, % Arrows on one end
	xlabel={$x$}, % Label the x-axis
	ylabel={$y$}, % Label the y-axis
}}

\pgfplotsset{mystyle/.style={no marks, line width=0.5pt, <->}} % Style of the graph

\title{Introduction to \LaTeX{}}
\subtitle{A \LaTeX{} Presentation}
\author{Vincent Macri\\JJ Marr}
\date{2017}

\usetheme{Rochester}
\usecolortheme{orchid}
\beamertemplatenavigationsymbolsempty
\usefonttheme[onlymath]{serif}
\setbeamercolor{background canvas}{bg=}
\logo{\color{logocolor}\rmfamily\normalsize\LaTeX{}}

\AtBeginSection[] {%
	\begin{frame}
	\frametitle{Lesson Progress}
	\tableofcontents[currentsection]
	\end{frame}
}

\begin{document}
	\frame{\titlepage}
	\section{Introduction}
	\begin{frame}
		\frametitle{\secname}
		\framesubtitle{Welcome to \LaTeX{}}
		\begin{itemize}[<+->]
			\item Lamport \TeX{}
			\item Written as \LaTeX{} or LaTeX.
			\item Pronounced lay-tek
		\end{itemize}
	\end{frame}
	\begin{frame}
		\frametitle{\secname}
		\framesubtitle{Why we use \LaTeX{}}
		\LaTeX{} is a very useful skill if you plan on doing:
		\pause
		\begin{itemize}
			\item Math
			\item Computer science
			\item Engineering
			\item Physics
			\item Chemistry
			\item Publishing
		\end{itemize}
		\pause
		Since \LaTeX{} is so versatile, once you get used to it, you will find that you can do almost anything in it.
	\end{frame}
	\begin{frame}
		\frametitle{\secname}
		\framesubtitle{Learning \LaTeX{}}
		The hardest part about learning \LaTeX{} is getting started.
		\pause

		So congratulations! You're done the hard part!
		\pause

		When using \LaTeX{}, it is important to remember that it makes correct design choices 99\% of the time. Accept what it does, as it does it for a reason.
		\pause

		For example, \LaTeX{} uses very large margins by default. It may look odd, but it actually sets the margins to make lines the size that studies have shown is best for readability (60--70 characters). This is easy to change of course.
	\end{frame}
	\begin{frame}
		\frametitle{\secname}
		\framesubtitle{Basic syntax}
		\begin{description}[<+->]
			\item[Spacing] Spaces in \LaTeX{} don't do much. \LaTeX{} treats 20 spaces the same as it treats 1 space.
			\item[Inline math] \lstinline{\$1 + 2 = 3\$} $\rightarrow 1 + 2 = 3$ In math mode, you can use 0 or 20 spaces, \LaTeX{} will output the same thing. Do whatever you find easiest to write and edit.
			\item[Comments] Comment \LaTeX{} code with the \% symbol. Write \textbackslash\% to use the \% symbol.
			\item[Commands] \LaTeX{} uses commands to perform quick actions. They are formatted as \textbackslash command\{parameter\}. For some commands, the \{\} and parameter are optional.
			\item[Environments] \LaTeX{} uses environments for things more advanced than commands. They start with \textbackslash begin\{environment\} and end with \textbackslash end\{environment\}
		\end{description}
	\end{frame}
	\begin{frame}
		\frametitle{\secname}
		\framesubtitle{Installing}
		Installing \LaTeX{} is easy, but there are lots of options for installing it.
		
		Listed below are the easiest solutions for the main desktop operating systems. Some of these work on all operating systems, but certain ones are easier with certain operating systems.
		\begin{description}
			\item[GNU/Linux and BSD] Install \TeX{} Live from your distro's repo
			\item[macOS] Mac\TeX{} is supposed to be good
			\item[Windows] MikTeX is a GUI program for editing and compiling \LaTeX{}
		\end{description}
		If you have difficulty installing \LaTeX{}, you can try an online editor such as overleaf.com

		You can also check \url{https://www.latex-project.org/get/} if you are not satisfied with any of these options.
	\end{frame}
	\begin{frame}
		\frametitle{\secname}
		\framesubtitle{Let's start}
		In this presentation, we will show how to write a \LaTeX{} document that proves the quadratic formula.
	\end{frame}
	\section{Preamble}
	\begin{frame}
		\frametitle{\secname}
		\framesubtitle{The beginning}
		Our document is split into two sections.

		The \alert{preamble} and the \alert{main document}.
	\end{frame}
	\begin{frame}
		\frametitle{\secname}
		\framesubtitle{The preamble}
		Your preamble is where you tell \LaTeX{} about what is going to be in your document.

		Information is the preamble includes information such as:
		\begin{itemize}[<+->]
			\item Document class
			\item Packages
			\item Document title
			\item Author
			\item Date (defaults to date when document was compiled)
		\end{itemize}
	\end{frame}
	\begin{frame}
		\frametitle{\secname}
		\framesubtitle{Example}
		Let's code our preamble.
		\lstinputlisting[linerange=1-12]{QuadraticProof.tex}
	\end{frame}
	\section{Setup}
	\begin{frame}[fragile]
		\frametitle{\secname}
		\framesubtitle{Subtitle}
		\LaTeX{} uses \alert{sections} to split up parts of the document.

		In school assignments, you will usually have each question be its own section.
		\begin{example}
			\begin{lstlisting}[numbers=none, gobble=16]
				\section{Prove that $1 + 1 = 2$}
			\end{lstlisting}
		\end{example}
	\end{frame}
	\begin{frame}[fragile]
		\frametitle{\secname}
		\framesubtitle{Title page and table of contents}
		The \lstinline{\maketitle} command is used to generate a title page.
		\begin{example}
			\begin{lstlisting}[numbers=none, gobble=16]
				\maketitle
			\end{lstlisting}
		\end{example}
		
		The \lstinline{\tableofcontents} command will generate a table of contents out of your sections.
		\begin{example}
			\begin{lstlisting}[numbers=none, gobble=16]
				\tableofcontents
			\end{lstlisting}
		\end{example}
		If your items aren't showing up in your table of contents, try recompiling.
	\end{frame}
	\section{Math}
	\begin{frame}[fragile]
		\frametitle{\secname}
		\framesubtitle{Doing math}
		When doing math, we frequently plug values into formula and solve for variables.

		The \alert{align*} environment is useful for this.
		\begin{example}
			\begin{lstlisting}[numbers=none, gobble=16]
				\begin{align*}
					a^2 + b^2 &= c^2\\
					a^2 &= c^2 - b^2
				\end{align*}
			\end{lstlisting}
		\end{example}
		\Large
		\begin{align*}
			a^2 + b^2 &= c^2\\
			a^2 &= c^2 - b^2
		\end{align*}
	\end{frame}
	\begin{frame}[fragile]
		\frametitle{\secname}
		\framesubtitle{Math symbols}
		Here is a short list of useful \LaTeX{} math symbols:
		\begin{center}
			\begin{columns}
				\begin{column}{0.5\textwidth}
					\begin{description}
						\item[$\therefore$] \lstinline{$\therefore$}
						\item[$\because$] \lstinline{$\because$}
						\item[$\approx$] \lstinline{$\approx$}
						\item[$\rightarrow$] \lstinline{$\rightarrow$}
						\item[$\neq$] \lstinline{$\neq$}
						\item[$\pm$] \lstinline{$\pm$}
						\item[$\times$] \lstinline{$\times$}
						\item[$\cdot$] \lstinline{$\cdot$}
						\item[$\infty$] \lstinline{$\infty$}
					\end{description}
				\end{column}
				\begin{column}{0.5\textwidth}
					\begin{description}
						\item[$\Delta$] \lstinline{$\Delta$}
						\item[$\pi$] \lstinline{$\pi$}
						\item[$\tau$] \lstinline{$\tau$}
						\item[$\theta$] \lstinline{$\theta$}
						\item[$\angle$] \lstinline{$\angle$}
						\item[$\cup$] \lstinline{$\cup$}
						\item[$\cap$] \lstinline{$\cap$}
						\item[$\geq$] \lstinline{$\ge$} or \lstinline{$\geq$}
						\item[$\leq$] \lstinline{$\le$} or \lstinline{$\leq$}
					\end{description}
				\end{column}
			\end{columns}
		\end{center}
		You can search the Internet for other commands.

		You can also try \url{http://detexify.kirelabs.org}, a site that lets you draw the symbol you want and then lists \LaTeX{} commands for what it thinks you're looking for.
	\end{frame}
	\begin{frame}[fragile]
		\frametitle{\secname}
		\framesubtitle{Math notation}
		\begin{columns}
			\begin{column}{0.6\textwidth}
				\begin{example}[Powers]
					\begin{lstlisting}[numbers=none, gobble=24]
						$a^{b}$
					\end{lstlisting}
				\end{example}
			\end{column}
			\begin{column}{0.4\textwidth}
				\Huge
				$a^{b}$
			\end{column}
		\end{columns}
		\begin{columns}
			\begin{column}{0.6\textwidth}
				\begin{example}[Subscript]
					\begin{lstlisting}[numbers=none, gobble=24]
						$x_{1}$
					\end{lstlisting}
				\end{example}
			\end{column}
			\begin{column}{0.4\textwidth}
				\Huge
				$x_{1}$
			\end{column}
		\end{columns}
		\begin{columns}
			\begin{column}{0.6\textwidth}
				\begin{example}[Element Of]
					\begin{lstlisting}[numbers=none, gobble=24]
						$x \in \mathbb{R}$
					\end{lstlisting}
				\end{example}
			\end{column}
			\begin{column}{0.4\textwidth}
				\Huge
				$x \in \mathbb{R}$
			\end{column}
		\end{columns}
	\end{frame}
	\begin{frame}[fragile]
		\frametitle{\secname}
		\framesubtitle{Math commands}
		\begin{columns}
			\begin{column}{0.6\textwidth}
				\begin{example}[Fractions]
					\begin{lstlisting}[numbers=none, gobble=24]
						$\frac{a}{b}$
					\end{lstlisting}
				\end{example}
			\end{column}
			\begin{column}{0.4\textwidth}
				\Huge
				$\frac{a}{b}$
			\end{column}
		\end{columns}
		\begin{columns}
			\begin{column}{0.6\textwidth}
				\begin{example}[Square Roots]
					\begin{lstlisting}[numbers=none, gobble=24]
						$\sqrt{a}$
					\end{lstlisting}
				\end{example}
			\end{column}
			\begin{column}{0.4\textwidth}
				\Huge
				$\sqrt{a}$
			\end{column}
		\end{columns}
		\begin{columns}
			\begin{column}{0.6\textwidth}
				\begin{example}[Nth Roots]
					\begin{lstlisting}[numbers=none, gobble=24]
						$\sqrt[n]{a}$
					\end{lstlisting}
				\end{example}
			\end{column}
			\begin{column}{0.4\textwidth}
				\Huge
				$\sqrt[n]{a}$
			\end{column}
		\end{columns}
	\end{frame}
	\begin{frame}[fragile]
		\frametitle{\secname}
		\framesubtitle{Quadratic Formula Example}
		\begin{example}
			\begin{lstlisting}[numbers=none, gobble=16]
				$x = \frac{-b \pm \sqrt{b^2 - 4ac}}{2a}$
			\end{lstlisting}
		\end{example}
		\begin{center}
			\Huge
			$x = \frac{-b \pm \sqrt{b^2 - 4ac}}{2a}$
		\end{center}
	\end{frame}
	\section{The Document}
	\begin{frame}[fragile]
		\frametitle{\secname}
		\framesubtitle{Let's actually start}
		Time to actually start on our document contents!

		Let's explain how we prove the quadratic formula using our words before we derive it.
		\begin{example}
			\begin{lstlisting}[numbers=none, gobble=16]
				\section{Prove the quadratic formula}
					We will prove this by deriving it from $ax^2 + bx + c = 0$.

					This is fairly trivial.
			\end{lstlisting}
		\end{example}
	\end{frame}
	\section{Graphing}
	\begin{frame}[fragile]
		\frametitle{\secname}
		\framesubtitle{PGFPlots and TikZ}
		To graph in \LaTeX{}, we use the package PGFPlots, which uses another package called TikZ.

		Graphing with these is very easy, but can also be verbose.

		In order to use PGFPlots, we must use its package in our preamble:
		\begin{lstlisting}[numbers=none, gobble=12]
			\usepackage{pgfplots}
		\end{lstlisting}
	\end{frame}
	\begin{frame}[fragile]
		\frametitle{\secname}
		\framesubtitle{Graph Options}
		In your preamble:
		\begin{lstlisting}[gobble=12]
			\pgfplotsset{every axis/.append style={
				axis x line=middle, % Make the x-axis position be the middle of the graph
				axis y line=middle, % Make the y-axis position be the middle of the graph
				axis line style={->}, % Arrows on one end
				xlabel={$x$}, % Label the x-axis
				ylabel={$y$}, % Label the y-axis
			}}

			\pgfplotsset{mystyle/.style={no marks, line width=0.5pt, <->}} % Style of the graph
		\end{lstlisting}
	\end{frame}
	\begin{frame}[fragile]
		\frametitle{\secname}
		\framesubtitle{Graph Options}
		\begin{lstlisting}[gobble=12]
			\begin{tikzpicture}
			\begin{axis}[
				width=\textwidth, % Make the graph take up the width of the page
				height=3in, % Make the graph 3 inches tall
				xmin=-5,xmax=5, % Make the x-axis go from -5 to 5
				ymin=-2.5,ymax=8, % Make the y-axis go from -2.5 to 8
				]
				\addplot[mystyle] [domain=-2:4, samples=50]{(x-1)^2 - 2} node at (axis cs:-3, 1){$f(x)=(x-1)^2 - 2$};
			\end{axis}
			\end{tikzpicture}
		\end{lstlisting}
	\end{frame}
	\begin{frame}
		\frametitle{\secname}
		\framesubtitle{The Graph}
		\begin{tikzpicture}
			\begin{axis}[
				width=\textwidth, % Make the graph take up the width of the page
				height=3in, % Make the graph 3 inches tall
				xmin=-5,xmax=5, % Make the x-axis go from -5 to 5
				ymin=-2.5,ymax=8, % Make the y-axis go from -2.5 to 8
				]
				\addplot[mystyle] [domain=-2:4, samples=50]{(x-1)^2 - 2} node at (axis cs:-3, 1){$f(x)=(x-1)^2 - 2$};
			\end{axis}
		\end{tikzpicture}
	\end{frame}
	\begin{frame}[fragile]
		\frametitle{\secname}
		\framesubtitle{Sinusoidal Curves}
		\begin{lstlisting}[gobble=12]
			\begin{tikzpicture}
				\begin{axis}[
					width=\textwidth, % Make the graph take up the width of the page
					height=3in, % Make the graph 3 inches tall
					xmin=0,xmax=2.25*pi, % Make the x-axis go from 0 to 2pi
					ymin=-1.5,ymax=1.5, % Make the y-axis go from -1.5 to 1.5
					xtick = {0, 1.5708, ..., 6.2832}, % Marked positions on the x-axis
					xticklabels={0$\pi$, 0.5$\pi$,...$\pi$,2$\pi$}, % How to label those positions.
					]
					\addplot[mystyle, ->] [domain=0:2.125*pi, samples=100]{sin(deg(x))};
				\end{axis}
			\end{tikzpicture}
		\end{lstlisting}
	\end{frame}
	\begin{frame}
		\frametitle{\secname}
		\framesubtitle{Sin Graph}
		\begin{tikzpicture}
			\begin{axis}[
					width=\textwidth, % Make the graph take up the width of the page
					height=3in, % Make the graph 3 inches tall
					xmin=0,xmax=2.25*pi, % Make the x-axis go from 0 to 2pi
					ymin=-1.5,ymax=1.5, % Make the y-axis go from -1.5 to 1.5
					xtick = {0, 1.5708, ..., 6.2832}, % Marked positions on the x-axis
					xticklabels={0$\pi$, 0.5$\pi$,...$\pi$,2$\pi$}, % How to label those positions.
				]
				\addplot[mystyle, ->] [domain=0:2.125*pi, samples=100]{sin(deg(x))};
			\end{axis}
		\end{tikzpicture}
	\end{frame}
	\section{The Quadratic Proof}
	\begin{frame}
		\frametitle{\secname}
		\framesubtitle{The preamble}
		\lstinputlisting[linerange=1-12]{QuadraticProof.tex}
	\end{frame}
	\begin{frame}[shrink]
		\frametitle{\secname}
		\framesubtitle{The content}
		\lstinputlisting[firstline=14, firstnumber=14]{QuadraticProof.tex}
	\end{frame}
	\section{Graphics}
	\begin{frame}[fragile]
		\frametitle{\secname}
		\framesubtitle{Look at these graphics}
		\LaTeX{} lets you include images with the \lstinline{graphicx} package.
		\begin{example}
			\begin{lstlisting}[gobble=12]
				\begin{center}
					\includegraphics[height=0.5in]{PolynomialSuperheroes} % The actual file is named PolynomialSuperheroes.eps
				\end{center}
			\end{lstlisting}
		\end{example}
		\begin{center}
			\includegraphics[height=0.5in]{PolynomialSuperheroes}
		\end{center}
		The \lstinline{\includegraphics} command does need the file extension. It will pick the best format if multiple formats are available. It supports: JPG, PNG, PDF, and (usually) EPS.
	\end{frame}
	\begin{frame}[fragile]
		\frametitle{\secname}
		\framesubtitle{Including entire pages}
		If you don't want to bother with using \LaTeX{} to graph, you can scan an image, and use the \lstinline{pdfpages} package to include the scan in your \LaTeX{} document.
		\begin{example}
			\begin{lstlisting}[numbers=none, gobble=16]
				\includepdf[pages={1-2}]{QuadraticProof.pdf} % Include JJ's quadratic proof.
			\end{lstlisting}
		\end{example}
		The next two slides will contain JJ's quadratic proof.
	\end{frame}
	\includepdf[pages={1-2}]{QuadraticProof.pdf} % Include JJ's quadratic proof.
	\section{Questions and Answers}
	\begin{frame}
		\frametitle{\secname}
		\framesubtitle{Question for me?}
		\begin{center}
			{\Huge Questions}
		\end{center}
	\end{frame}
\end{document}
